% !TEX root = ../../tree_notes.tex
% =====================================================================
%  Chapter 8 opener -- "Chapter at a Glance", reading paths + ladder
%  staircase, the two running examples (the degenerate chain C and its
%  balanced twin; the zig-zag Z that forces a double rotation), the
%  "who balances, how" comparison, roadmap, recognition guide, notation.
%  NO numbered boxes here.
% =====================================================================

\begin{chaptertheorybox}{Chapter at a Glance}
Chapter~7 (\chref{c05}) spent one invariant -- $\text{left}<\text{node}<\text{right}$ -- to turn a
binary tree into a search index, and closed on a debt: every routine costs $\bigO{h}$, and $h$ is
\emph{not yours to control}. Feed the keys in sorted order and the BST rots into a linked list, so
``search a million times'' silently becomes ``scan a million times.'' This chapter pays that debt by
\emph{construction}: it keeps the height logarithmic, forever, no matter what order the keys arrive
in. There are exactly four ways to do it, and this chapter teaches all four dialects of one sentence
-- \emph{restore the shape without ever breaking the inorder line}:
\begin{itemize}[nosep]
  \item \textbf{The rotation is the only legal move.} A single local pointer surgery re-weights a
        subtree while leaving the sorted inorder sequence \emph{identical}. Every balanced tree, from
        AVL to red-black to splay, is a policy for \emph{when} to rotate and \emph{how much} to care.
        \hfill \lv{2}--\lv{3}
  \item \textbf{Three philosophies of staying short.} \emph{Deterministic rebalancing} (AVL) enforces
        a strict height rule and rotates on every edit; \emph{randomisation} (treap) trades case work
        for a coin flip and wins $\bigO{\log n}$ in expectation; \emph{colour invariants} (red-black)
        relax the rule just enough that $\le 2$ recolours/rotisons suffice -- which is \emph{why} your
        \code{std::set} is fast. \hfill \lv{3}--\lv{5}
  \item \textbf{A height you control unlocks augmentation.} Once $h=\bigO{\log n}$ by construction, a
        size-augmented or implicit-keyed tree answers rank, select, and range questions on a
        \emph{dynamic sequence} -- the promised, guaranteed version of the fragile trick the Chapter~7
        Summit (\probref{c05-summit}) could only limp through. \hfill \lv{4}--\lv{5}
\end{itemize}

\textbf{You will learn to}
\begin{itemize}[nosep]
  \item perform and reason about the four rotation shapes (LL, RR, LR, RL) as the unique
        inorder-preserving fixes for a height violation; \hfill\lv{2}--\lv{3}
  \item insert and delete in an AVL tree, tracking balance factors up the access path and stopping at
        the first restored ancestor; \hfill\lv{3}--\lv{4}
  \item build a treap by assigning random priorities, and wield \code{split}/\code{merge} as the
        universal toolkit for insert, erase, rank, and range operations; \hfill\lv{4}
  \item explain, from the red-black invariants, why \code{std::set}/\code{std::map} guarantee
        $\bigO{\log n}$ -- the promise Chapter~7 borrowed and now repays; \hfill\lv{3}--\lv{4}
  \item recognise an \emph{implicit-key} treap as a dynamic array (insert/delete by index, reverse a
        range) and a splay tree as an amortised, locality-friendly container. \hfill\lv{4}--\lv{5}
\end{itemize}

\textbf{Prerequisites.} Chapter~7 (\chref{c05}) in full -- the BST descent, insertion, three-case
deletion, and the $\bigO{h}$-is-the-whole-cost warning are assumed, never re-derived. The inorder
walk is a black box from Chapter~3 (\chref{c03}); here we only \emph{lean on} the fact that a
rotation leaves it unchanged.
\end{chaptertheorybox}

% ---------------- Station 1b: reading paths + ladder position ----------------
\begin{chaptertheorybox}{How to Read This Chapter (two paths up the ladder)}
Level badges $\lv1$--$\lv6$: $\lv2$ is ``do the rotation by hand,'' $\lv3$ is a solid interview
medium, $\lv4$ is implementing a whole balanced structure cold, and $\lv5$ is an LC Hard where the
naive BST's height is the enemy the structure exists to slay.

\textbf{Beginner path ($\lv2$--$\lv3$).} Read \secref{c06-rotations} end to end -- it fixes \emph{why}
height is cost and gives you the four rotations as the one legal repair. Then \secref{c06-avl} through
insertion and the four rebalance cases on the running examples $C$ and $Z$. Take every
$\lv2$--$\lv{3}$ solved problem. You may \emph{skip on a first pass}: treap \code{split}/\code{merge}
and implicit keys, the red-black invariant proof sketch, splay amortisation, and the Summit. By the
end you can keep a search tree of a million keys at depth about $20$ no matter how they arrive.

\textbf{Full path.} Everything, in order -- it ends at the scaled Summit (LC Hard 493 \emph{Reverse
Pairs}: count $i<j$ with $\mathit{nums}[i]>2\cdot\mathit{nums}[j]$, where the guaranteed-$\bigO{\log
n}$ rank query of a balanced, augmented tree is the honest upgrade over Chapter~7's fragile descent,
and merge/Fenwick are the array-indexed rivals). The reduction gallery shows that ``AVL,'' ``treap,''
``red-black'' and ``splay'' are the same rotation-policy question with different answers, and that an
implicit treap is just a balanced BST whose keys are \emph{positions}. The forward rope is Chapter~9
(\chref{c26}, Cartesian trees and the monotonic stack) for the array$\to$tree bridge this chapter
deliberately leaves out, and Chapter~25 (\chref{c23}) for the full splay/link-cut machinery.
\end{chaptertheorybox}

\begin{figure}[H]
\centering
\begin{tikzpicture}[diagram, x=9mm, y=7mm]
  \foreach \st/\lab in {1/lv1,2/lv2,3/lv3,4/lv4,5/lv5,6/lv6} {
    \fill[\lab!22, draw=\lab, line width=0.8pt] (\st-1,0) rectangle (\st,\st*0.52);
    \node[font=\scriptsize\bfseries, text=\lab] at (\st-0.5,\st*0.52-0.42) {L\st};}
  \node[annot, anchor=south west, font=\scriptsize]
    at (0.05,0.62) {8.1 rotations};
  \node[annot, anchor=south west, font=\scriptsize]
    at (1.05,1.14) {8.1/8.2 the four rebalance cases};
  \node[annot, anchor=south west, font=\scriptsize]
    at (2.05,1.66) {8.2 AVL insert/delete cold};
  \node[annot, anchor=south west, font=\scriptsize]
    at (3.05,2.18) {8.3 treaps; split/merge; implicit keys};
  \node[annot, anchor=south west, font=\scriptsize, text=lv5]
    at (4.05,2.70) {8.4 red-black proof; 8.5 splay amortisation};
  \node[annot, anchor=south west, font=\scriptsize, text=lv5]
    at (5.05,3.22) {scaled summit (493)};
\end{tikzpicture}
\caption*{Chapter~8 on the six-step ladder: doing a rotation by hand ($\lv2$) climbs to a full AVL
insert/delete ($\lv3$); treaps and their split/merge engine and the red-black guarantee sit at
$\lv3$--$\lv4$; implicit-sequence treaps and splay amortisation reach $\lv4$; the Summit (493) makes
the controlled height pay for itself against an LC Hard, at $\lv5$. Beginner path $=$ steps 1--3.}
\end{figure}

% ---------------- The running examples ----------------
\subsection*{The Running Examples: A Chain That Must Be Saved, and a Case That Needs Two Turns}

This chapter has \emph{two} running examples, used over and over. The first is the villain: insert
the keys $1,2,\dots,7$ \emph{in increasing order} into an ordinary BST of Chapter~7 and watch it
collapse into a right-going chain of height $6$ -- search degrades to a scan, exactly the cliffhanger
that chapter ended on. The second is the interesting one: a violation that no \emph{single} rotation
can fix, which forces the LR double rotation. Every topic revisits them -- the rotation ($8.1$), AVL
rebalancing ($8.2$), treap priorities and split/merge ($8.3$), red-black recolouring ($8.4$), and the
splay access ($8.5$).

\begin{figure}[H]
\centering
\begin{tikzpicture}[diagram, scale=0.82, transform shape]
  % left: the chain C (ascending insert), height 6
  \node[tleaf]  (c1) at (0,0)     {$1$};
  \node[tleaf]  (c2) at (1.15,-0.9) {$2$};
  \node[tleaf]  (c3) at (2.3,-1.8)  {$3$};
  \node[tleaf]  (c4) at (3.45,-2.7) {$4$};
  \node[tleaf]  (c5) at (4.6,-3.6)  {$5$};
  \node[tleaf]  (c6) at (5.75,-4.5) {$6$};
  \node[tnode]  (c7) at (6.9,-5.4)  {$7$};
  \foreach \a/\b in {c1/c2,c2/c3,c3/c4,c4/c5,c5/c6,c6/c7}
    \draw[tedgehl] (\a) -- (\b);
  \node[annot, align=left, font=\scriptsize] at (1.4,0.9)
    {\textbf{Chain $C$}: insert $1,2,\dots,7$ in order\\ height $=6=n-1$\\ every descent is now
     $\bigO{n}$ -- the rot};
  % right: the balanced twin
  \node[troot]  (r4) at (12,0)   {$4$};
  \node[tnode]  (r2) at (10.4,-1.3) {$2$};
  \node[tnode]  (r6) at (13.6,-1.3) {$6$};
  \node[tleaf]  (r1) at (9.6,-2.6)  {$1$};
  \node[tleaf]  (r3) at (11.2,-2.6) {$3$};
  \node[tleaf]  (r5) at (12.8,-2.6) {$5$};
  \node[tleaf]  (r7) at (14.4,-2.6) {$7$};
  \foreach \a/\b in {r4/r2,r4/r6,r2/r1,r2/r3,r6/r5,r6/r7}
    \draw[tedgehl] (\a) -- (\b);
  \node[annot, align=left, font=\scriptsize] at (12,1.1)
    {\textbf{Balanced twin}: same seven keys\\ height $=2$\\ a rotation policy \emph{keeps} it here};
  \draw[flow, double] (7.6,-2.6) -- (9.2,-2.6) node[midway, above, font=\scriptsize] {balance};
\end{tikzpicture}
\caption{The villain, $C$: the \emph{same} set of keys stored as a degenerate chain (left, height $6$)
and as a height-balanced tree (right, height $2$). Nothing changes about the \emph{contents} or the
inorder line -- both read $1,2,3,4,5,6,7$ -- only the \emph{shape}. Balance is the art of moving from
left to right and \emph{staying} there, on every future edit, using the one legal move: the rotation.}\label{fig:c06-chain}
\end{figure}

\begin{figure}[H]
\centering
\begin{tikzpicture}[diagram, scale=0.86, transform shape]
  % Z before: 30 -> (10 -> (right 20 -> right 25)), 40
  \node[troot] (z30) at (0,0)     {$30$};
  \node[tnode] (z10) at (-1.6,-1.3) {$10$};
  \node[tnode] (z40) at (1.6,-1.3)  {$40$};
  \node[tnode] (z20) at (-0.5,-2.6) {$20$};
  \node[tleaf] (z25) at (0.4,-3.9)  {$25$};
  \draw (z30) -- (z10); \draw (z30) -- (z40);
  \draw (z10) -- (z20); \draw (z20) -- (z25);
  \node[annot, align=left, font=\scriptsize] at (-4.6,-2.6)
    {\textbf{$Z$, before}: insert $30,10,40,$\\ $20,25$. Node $30$ has\\ balance factor $+2$, its\\
     heavy child $10$ leans \emph{right}\\ $\Rightarrow$ the LR (zig-zag) case};
  % arrow
  \draw[flow, double] (3.0,-2.0) -- (4.6,-2.0) node[midway, above, font=\scriptsize, align=center] {rotate-left $10$,\\ then right $30$};
  % Z after: 20 -> 10,30 ; 30 -> 25,40
  \node[troot] (y20) at (7.2,-0.4) {$20$};
  \node[tleaf] (y10) at (5.8,-1.7) {$10$};
  \node[tnode] (y30) at (8.6,-1.7) {$30$};
  \node[tleaf] (y25) at (7.8,-3.0) {$25$};
  \node[tleaf] (y40) at (9.4,-3.0) {$40$};
  \draw (y20) -- (y10); \draw (y20) -- (y30);
  \draw (y30) -- (y25); \draw (y30) -- (y40);
  \node[annot, align=left, font=\scriptsize] at (7.6,0.7)
    {\textbf{$Z$, after}: height $2$, balanced\\ \emph{inorder unchanged}: $10,20,25,30,40$};
\end{tikzpicture}
\caption{The interesting case, $Z$: a Left-Right (zig-zag) violation. When the heavy child leans the
\emph{opposite} way, one rotation only converts the problem into a single-side case -- you need
\emph{two}: a rotation on the child first, then on the anchor. Both figures keep the identical sorted
inorder; this is the crux of \thmref{c06-rotation}.}\label{fig:c06-zigzag}
\end{figure}

\medskip
\noindent Keep the four \emph{strategies} below in mind as variations on one theme; each is proved or
used, never merely asserted, in the section named after it:
\begin{center}
\renewcommand{\arraystretch}{1.25}
\setlength{\tabcolsep}{5pt}
\begin{tabular}{@{}p{0.18\linewidth} p{0.44\linewidth} p{0.30\linewidth}@{}}
\toprule
\textbf{Scheme} & \textbf{How it keeps $h=\bigO{\log n}$} & \textbf{Where it lives} \\
\midrule
AVL          & strict rule: $\abs{\mathit{BF}}\le 1$ per node; rotate on every edit that breaks it & \secref{c06-avl} \\
Treap        & random priority per node; BST on keys \emph{and} heap on priorities is short \emph{in expectation} & \secref{c06-treaps} \\
Red-black    & colour invariants; longest path $\le 2\times$ shortest $\Rightarrow h\le 2\log_2(n{+}1)$ & \secref{c06-redblack} \\
Splay        & no stored rule: every access rotates to the root; amortised $\bigO{\log n}$ & \secref{c06-splay} \\
\bottomrule
\end{tabular}
\end{center}

% ---------------- Roadmap diagram ----------------
\begin{figure}[H]
\centering
\begin{tikzpicture}[diagram, node distance=3.5mm,
    top/.style={boxnode, fill=primaryblue!10, draw=primaryblue, text width=19mm,
                minimum height=10mm, font=\small},
    res/.style={annot, text width=19mm, anchor=north, font=\scriptsize}]
  \node[top] (t1) {\textbf{8.1} Why \& Rotations};
  \node[top, right=of t1] (t2) {\textbf{8.2} AVL};
  \node[top, right=of t2] (t3) {\textbf{8.3} Treaps};
  \node[top, right=of t3] (t4) {\textbf{8.4} Red-Black};
  \node[top, right=of t4] (t5) {\textbf{8.5} Splay};
  \foreach \a/\b in {t1/t2, t2/t3, t3/t4, t4/t5} \draw[flow] (\a) -- (\b);
  \node[res] at ([yshift=-2mm]t1.south) {the one legal move; four cases};
  \node[res] at ([yshift=-2mm]t2.south) {balance factors; $\le2$ turns fix};
  \node[res] at ([yshift=-2mm]t3.south) {random priorities; split/merge};
  \node[res] at ([yshift=-2mm]t4.south) {the std::set guarantee};
  \node[res] at ([yshift=-2mm]t5.south) {self-adjust; amortised};
\end{tikzpicture}
\caption*{The five topics are one question at increasing abstraction. 8.1 gives the \emph{atomic
move}; 8.2 uses it \emph{greedily and provably} (AVL); 8.2 hands 8.3's treap the same job via
\emph{chance} instead of case work, and the split/merge engine; 8.4 explains the \emph{relaxed} rule
the standard library actually ships; 8.5 drops the stored invariant entirely for an \emph{amortised}
self-adjusting tree that anticipates the dynamic-trees chapter.}
\end{figure}

% ---------------- Topic table ----------------
\begin{center}
\renewcommand{\arraystretch}{1.25}
\setlength{\tabcolsep}{4pt}
\begin{tabular}{@{}p{0.235\linewidth} p{0.365\linewidth} c p{0.23\linewidth}@{}}
\toprule
\textbf{Topic} & \textbf{Key results} & \textbf{Levels} & \textbf{Solved problem} \\
\midrule
\secref{c06-rotations}\ \ Why Balance? The Rotation
  & op cost $=h$; the rotation preserves inorder; LL/RR/LR/RL cases & $\lv2$--\lv{3}
  & \probref{c06-treecons} \\
\addlinespace
\secref{c06-avl}\ \ AVL Trees
  & balance factor; Fibonacci height bound; insert/delete with $\le2$ turns & $\lv3$--\lv{4}
  & \probref{c06-avl} \\
\addlinespace
\secref{c06-treaps}\ \ Treaps, Split/Merge, Implicit
  & BST+heap on random priorities; $\bigO{\log n}$ expected; split/merge kit & $\lv4$--\lv{5}
  & \probref{c06-rpmq} \\
\addlinespace
\secref{c06-redblack}\ \ Red-Black Trees (conceptual)
  & five colour invariants; $h\le 2\log_2(n{+}1)$ -- the std::set promise & $\lv3$--\lv{4}
  & (references \chref{c05}) \\
\addlinespace
\secref{c06-splay}\ \ Splay Trees (intro)
  & move-to-root; zig/zig-zig/zig-zag; amortised $\bigO{\log n}$ & $\lv4$--\lv{5}
  & (intro; full $\to$ \chref{c23}) \\
\bottomrule
\end{tabular}
\end{center}

% ---------------- Problem recognition guide ----------------
\subsection*{Problem Recognition Guide --- When You See \emph{this}, Think \emph{that}}

Balancing problems announce themselves when a \emph{plain BST or a hand-rolled ordered structure
must survive adversarial order}, or when you need a \emph{dynamic sequence} with range operations.
Match the phrasing; every row points at the section that builds the tool.

\begin{center}
\renewcommand{\arraystretch}{1.3}
\setlength{\tabcolsep}{4pt}
\begin{tabular}{@{}p{0.44\linewidth} p{0.34\linewidth} l@{}}
\toprule
\textbf{When the problem says\ldots} & \textbf{\textcolor{deepblue}{\textbf{Think\ldots}}} & \textbf{Section} \\
\midrule
``insert these keys; the tree must stay efficient'', keys arrive \emph{sorted}
  & plain BST $=$ a chain -- balance by construction & \secref{c06-rotations} \\
\addlinespace
``implement a self-balancing BST / AVL'', ``rebalance after insert''
  & track balance factors up the path; the four cases & \secref{c06-avl} \\
\addlinespace
``cut/paste a sequence, reverse a range, $k$th by position''
  & implicit-key treap; \code{split}/\code{merge} & \secref{c06-treaps} \\
\addlinespace
``rank / count-smaller / order statistic, online, worst-case fast''
  & size-augmented \emph{balanced} tree (or Fenwick) & \secref{c06-treaps}, Summit \\
\addlinespace
``why is \code{std::set} $\bigO{\log n}$?'' / a guarantee you must prove
  & red-black colour invariants & \secref{c06-redblack} \\
\addlinespace
``frequently accessed items should be cheap / no stored balance''
  & splay: move-to-root, amortised & \secref{c06-splay} \\
\bottomrule
\end{tabular}
\end{center}

\begin{center}
\renewcommand{\arraystretch}{1.25}
\setlength{\tabcolsep}{4pt}
\begin{tabular}{@{}p{0.27\linewidth} p{0.45\linewidth} p{0.20\linewidth}@{}}
\toprule
\textbf{Constraint you see} & \textbf{What it forces} & \textbf{Chapter-8 tool} \\
\midrule
keys arrive in (near-)sorted order & never trust a plain BST's height & \secref{c06-avl} \\
an ordered set/map needed under edits & reach for \code{set}/\code{map}, know it is red-black & \secref{c06-redblack} \\
a \emph{sequence} with range reverse/sum & implicit treap, keys are positions & \secref{c06-treaps} \\
rank/count queries, worst-case $\bigO{\log n}$ & augment with subtree size on a balanced tree & Summit (\probref{c06-summit}) \\
dense integer keys, array-index ranks & Fenwick / segment beat a pointer tree & \chref{c14}, \chref{c13} \\
\bottomrule
\end{tabular}
\end{center}

\begin{notebox}[Not this chapter?]
If you just need the \emph{ordered descent itself} -- search, insert, delete, kth-smallest, the
\code{std::set} API, interval management -- that is Chapter~7 (\chref{c05}); this chapter assumes all
of it and only asks \emph{how the height stays small}. If the tree must support \emph{link/cut of
whole subtrees} and path queries, the full splay and link-cut machinery is Chapter~25
(\chref{c23}) -- \secref{c06-splay} is only its prerequisite intro. If you want the
\emph{array$\to$tree bridge} -- Cartesian trees, the monotonic stack, largest rectangle, subarray
minimums, next-smaller links -- that is Chapter~9 (\chref{c26}), which this chapter deliberately
leaves out. And if keys are dense integers and you only need ranks by index, the Fenwick tree
(\chref{c14}) and segment tree (\chref{c13}) are faster engines for the very same sums. This chapter
owns \emph{keeping the ordered tree short by construction, and what that guaranteed height buys you}.
\end{notebox}

% ---------------- Notation ----------------
\subsection*{Notation Used in This Chapter}

\begin{center}
\renewcommand{\arraystretch}{1.2}
\setlength{\tabcolsep}{4pt}
\begin{tabular}{@{}p{0.20\linewidth} p{0.27\linewidth}@{\quad}p{0.17\linewidth} p{0.27\linewidth}@{}}
\toprule
\textbf{Symbol} & \textbf{Meaning} & \textbf{Symbol} & \textbf{Meaning} \\
\midrule
$h$                 & tree height; every descent costs $\bigO{h}$, balance forces $h=\bigO{\log n}$ & $\bal(v)$ & balance factor of $v$, $\height(L)-\height(R)$ \\
$C$                 & the degenerate chain (\figref{c06-chain})   & $Z$ & the zig-zag instance (\figref{c06-zigzag}) \\
$\mathit{p}(v)$     & random priority of $v$ in a treap (heap order) & $\mathit{size}(v)$ & node count of $v$'s subtree (augmentation) \\
$\mathrm{rot}_{L},\ \mathrm{rot}_{R}$ & left / right rotation about a node & $\mathrm{size}$-BF & implicit key $=$ left-subtree size \\
\code{set}, \code{map} & the \texttt{std} containers -- red-black under the hood (\secref{c06-redblack}) & $\lv{3}$ & difficulty badge, one of $\lv{1}$--$\lv{6}$ \\
\bottomrule
\end{tabular}
\end{center}
